\section{Objet de l'annexe}
\label{annexe-interfaces-objet}

Cette annexe développe les calculs associés aux deux premières
interfaces du Tome II.

\begin{itemize}
\item
  R1 transforme l'amplitude transmise par une perméabilité dynamique
  \(q\).
\item
  R2 transforme la temporalité de l'écriture : l'influence s'accumule
  dans une variable d'interface \(\chi\), puis entre dans la trace par
  quanta lors des Kairos.
\end{itemize}

L'objectif est de montrer précisément :

\begin{enumerate}
\item
  comment R1 transporte ou rend inaccessibles les seuils du modèle
  continu ;
\item
  comment construire l'application de retour de R2 ;
\item
  comment obtenir le seuil d'activité événementielle \(g_K\) ;
\item
  comment calculer la trace moyenne produite par les quanta ;
\item
  comment définir la stabilité hybride du repos et le seuil \(g_F\) ;
\item
  pourquoi les seuils \(g_K\), \(g_F\),
  \(g_{\mathrm{app}}^{R2}\) et \(g_H^{R2}\) doivent rester distincts.
\end{enumerate}

Les calculs sont effectués sur la réduction diagonale utilisée dans le
Tome II. Le système source est maintenu sur son puits positif, et

\[
w=\tanh(x_A)>0
\]

désigne l'intensité projetée par composante.

\section{R1 : perméabilité dynamique}

\subsection{Modèle}
\label{annexe-r1-modele}

La perméabilité \(q\) vérifie

\[
\dot q
=
\rho(1-q)
-
\eta qgw,
\]

avec

\[
\rho>0,
\qquad
\eta>0,
\qquad
g\geq0.
\]

Le terme \(\rho(1-q)\) rouvre l'interface en l'absence de trafic. Le
terme \(\eta qgw\) la ferme progressivement lorsque l'exposition est
soutenue.

L'écriture continue dans la trace n'est plus

\[
gw,
\]

mais

\[
q(t)gw.
\]

On suppose généralement

\[
0\leq q(0)\leq1.
\]

\subsection{Solution exacte de la perméabilité}
\label{annexe-r1-solution}

L'équation s'écrit

\[
\dot q
+
(\rho+\eta gw)q
=
\rho.
\]

Posons

\[
a_R(g)
=
\rho+\eta gw.
\]

La solution est

\[
\boxed{
q(t)
=
q_\infty(g)
+
\bigl(q_0-q_\infty(g)\bigr)
e^{-a_R(g)t},
}
\]

où

\[
\boxed{
q_\infty(g)
=
\frac{\rho}
{\rho+\eta gw}.
}
\]

L'équilibre \(q_\infty(g)\) est globalement exponentiellement stable,
avec un temps caractéristique

\[
\boxed{
t_R(g)
=
\frac{1}{\rho+\eta gw}.
}
\]

Si \(q_0\in[0,1]\), alors \(q(t)\in[0,1]\) pour tout \(t\geq0\).

En effet, le champ au bord vérifie

\[
\dot q\big|_{q=0}=\rho>0,
\]

et

\[
\dot q\big|_{q=1}=-\eta gw\leq0.
\]

L'intervalle \([0,1]\) est donc invariant.

\subsection{Couplage effectif}
\label{annexe-r1-couplage-effectif}

L'intensité transmise par l'interface est

\[
g_{\mathrm{eff}}(t;g)
=
gq(t).
\]

À long terme,

\[
\boxed{
g_{\mathrm{eff}}^\infty(g)
=
gq_\infty(g)
=
\frac{\rho g}
{\rho+\eta gw}.
}
\]

Cette fonction est strictement croissante :

\[
\frac{\mathrm d}{\mathrm dg}
g_{\mathrm{eff}}^\infty(g)
=
\frac{\rho^2}
{\left(\rho+\eta gw\right)^2}
>0,
\]

mais elle est bornée :

\[
\boxed{
\lim_{g\to\infty}
g_{\mathrm{eff}}^\infty(g)
=
\frac{\rho}{\eta w}.
}
\]

R1 ne supprime donc pas nécessairement l'influence. Elle lui impose un
plafond soutenu.

\section{Transport d'un seuil du modèle continu}
\label{annexe-r1-transport-seuil}

Soit

\[
g_\star^0>0
\]

un seuil du modèle continu R0. Par exemple, \(g_\star^0\) peut être le
seuil transcritique \(g_c^0\) ou le seuil de Hopf \(g_H^0\).

Le seuil quasi-statique de R1 est défini par

\[
g_{\mathrm{eff}}^\infty
\left(
g_\star^{R1}
\right)
=
g_\star^0.
\]

On résout

\[
\frac{\rho g_\star^{R1}}
{\rho+\eta wg_\star^{R1}}
=
g_\star^0.
\]

Après multiplication,

\[
\rho g_\star^{R1}
=
\rho g_\star^0
+
\eta wg_\star^0g_\star^{R1}.
\]

Ainsi,

\[
\left(
\rho-\eta wg_\star^0
\right)
g_\star^{R1}
=
\rho g_\star^0.
\]

Si

\[
\rho>\eta wg_\star^0,
\]

on obtient le seuil fini

\[
\boxed{
g_\star^{R1}
=
\frac{\rho g_\star^0}
{\rho-\eta wg_\star^0}.
}
\]

Si

\[
\rho\leq\eta wg_\star^0,
\]

alors

\[
\frac{\rho}{\eta w}
\leq
g_\star^0.
\]

Le plafond de R1 reste sous le seuil continu, ou ne l'atteint qu'à la
limite \(g\to\infty\). Aucun \(g\) fini ne permet donc de franchir ce
seuil en régime soutenu.

\begin{quote}
\resultheading{Proposition}{prop:E-transport-seuil-R1}{}
R1 transporte un seuil \(g_\star^0\) par la formule précédente lorsque
\[
\frac{\rho}{\eta w}>g_\star^0.
\]
Lorsque
\[
\frac{\rho}{\eta w}\leq g_\star^0,
\]
le seuil est inaccessible par une augmentation finie de \(g\).
\end{quote}

\section{Protection absolue et protection de la résistance}
\label{annexe-r1-protection}

Deux cas sont particulièrement importants.

\subsection{Protection du repos}

Si

\[
\frac{\rho}{\eta w}
\leq
g_c^0,
\]

le couplage soutenu ne peut jamais atteindre le seuil qui déstabilise
le repos. L'interface protège alors le système receveur contre tout
allumage quasi-statique, quelle que soit l'amplitude extérieure brute.

\subsection{Protection de la branche résistante}

Si

\[
g_c^0
<
\frac{\rho}{\eta w}
\leq
g_H^0,
\]

le repos peut perdre sa stabilité, mais le seuil de Hopf de la branche
résistante reste inaccessible.

La résistance peut donc apparaître sans pouvoir être rompue par une
simple augmentation de l'amplitude soutenue.

Pour les paramètres étudiés,

\[
\rho=0{,}05,
\qquad
\eta=1,
\]

et

\[
\frac{\rho}{\eta w}
=
0{,}0510.
\]

Comme

\[
g_c^0=0{,}0327
<
0{,}0510
<
g_H^0=0{,}1285,
\]

R1 autorise l'apparition de la résistance, mais protège
structurellement cette branche contre la rupture de Hopf du modèle
continu.

\section{Fenêtre initiale d'ouverture}
\label{annexe-r1-fenetre-initiale}

La formule de transport précédente porte sur le régime soutenu. Elle
ne décrit pas exactement un protocole démarrant avec une interface
ouverte.

Si

\[
q_0>q_\infty(g),
\]

alors

\[
g_{\mathrm{eff}}(t;g)
>
g_{\mathrm{eff}}^\infty(g)
\]

pour tout temps fini.

L'excès transmis au-dessus du régime asymptotique vaut

\[
g_{\mathrm{eff}}(t;g)
-
g_{\mathrm{eff}}^\infty(g)
=
g
\bigl(
q_0-q_\infty(g)
\bigr)
e^{-a_R(g)t}.
\]

Son intégrale totale est

\[
\boxed{
\int_0^\infty
\left[
g_{\mathrm{eff}}(t;g)
-
g_{\mathrm{eff}}^\infty(g)
\right]
\,\mathrm dt
=
\frac{
g\bigl(q_0-q_\infty(g)\bigr)
}{
\rho+\eta gw
}.
}
\]

Cette quantité mesure la fenêtre initiale d'ouverture.

Avec \(q_0=1\), le seuil quasi-statique calculé pour
\(g_c^{R1}\) vaut

\[
0{,}0913,
\]

alors que le seuil observé sous mise en contact brutale vaut

\[
0{,}0863.
\]

Cette différence ne contredit pas la formule analytique : le système
profite du surcroît initial de perméabilité avant la fermeture de
l'interface.

\section{Absence d'hystérésis dans R1 minimal}
\label{annexe-r1-absence-hysteresis}

Pour chaque valeur fixée de \(g\), la dynamique de \(q\) possède un
équilibre unique \(q_\infty(g)\), globalement stable.

Il n'existe donc ni branche ouverte ni branche fermée coexistantes à
la même valeur de \(g\). La loi R1 minimale ne produit pas
d'hystérésis propre.

Une hystérésis demanderait un mécanisme supplémentaire, par exemple :

\begin{itemize}
\item
  deux seuils distincts d'ouverture et de fermeture ;
\item
  une usure cumulative ;
\item
  une variable lente modifiant durablement \(\rho\) ou \(\eta\).
\end{itemize}

\section{Effet sur la durée d'appropriation}
\label{annexe-r1-duree}

R1 réduit l'amplitude transmise après la fenêtre initiale. La charge de
trace monte donc plus lentement et la durée minimale d'exposition
augmente.

Pour le protocole de référence :

\[
\begin{array}{c|cc}
g
&
T^\ast\text{ sous R0}
&
T^\ast\text{ sous R1}
\\
\hline
0{,}10&24{,}8&88{,}9\\
0{,}20&15{,}0&32{,}1\\
0{,}30&12{,}3&26{,}2
\end{array}
\]

Ces valeurs sont numériques. La partie démontrée ici est la loi de
filtrage \(g\mapsto g_{\mathrm{eff}}(t;g)\), non une formule fermée
pour \(T^\ast\).

\section{R2 : écriture événementielle}

\subsection{Modèle réduit au repos}
\label{annexe-r2-modele}

R2 ne laisse plus l'influence écrire directement dans la trace. Elle
s'accumule d'abord dans une variable d'interface \(\chi\).

Sur la diagonale positive, chaque composante de \(\chi\) vérifie

\[
\dot\chi
=
gw-\lambda_\chi\chi.
\]

Comme le vecteur complet vaut \((\chi,\chi)\),

\[
\|\boldsymbol\chi\|
=
\sqrt2\,\chi.
\]

Au repos exact du système receveur,

\[
x=0,
\qquad
v=0,
\]

la pression vérifie

\[
\dot p
=
\kappa_\chi\sqrt2\,\chi
-
\kappa_3p.
\]

Lors d'un Kairos,

\[
\theta^+
=
\theta^-
+
\kappa_Q\chi^-,
\]

\[
\chi^+
=
(1-r_\chi)\chi^-,
\]

\[
p^+
=
r_P\Theta.
\]

La variable \(\theta\) désigne ici une composante de la trace
diagonale

\[
\tau=(\theta,\theta).
\]

Entre les événements, au repos exact,

\[
\dot\theta=-\beta\theta.
\]

Les paramètres de référence sont

\[
\lambda_\chi=0{,}2,
\qquad
\kappa_Q=0{,}6,
\qquad
r_\chi=0{,}8,
\qquad
\kappa_\chi=2{,}5.
\]
Aux paramètres de référence, \(\kappa_3=0{,}5\neq
\lambda_\chi=0{,}2\) : la forme générique de la solution ci-dessous
est donc celle effectivement utilisée.

Cette règle réalise une écriture événementielle. Elle n'est pas encore
conditionnelle à une compatibilité entre le contenu reçu et l'état de
l'interface.

\subsection{Solution de \texorpdfstring{$\chi$}{chi} entre deux événements}
\label{annexe-r2-solution-chi}

Posons

\[
\chi_\infty(g)
=
\frac{gw}{\lambda_\chi}.
\]

À partir d'une valeur post-événement \(\chi_0\),

\[
\boxed{
\chi(t)
=
\chi_\infty
+
\left(
\chi_0-\chi_\infty
\right)
e^{-\lambda_\chi t}.
}
\]

La variable \(\chi\) tend donc vers la charge d'interface
\(\chi_\infty\) tant que le contact reste actif.

\subsection{Solution de la pression}
\label{annexe-r2-solution-pression}

Posons

\[
c_\chi
=
\kappa_\chi\sqrt2.
\]

La pression satisfait

\[
\dot p+\kappa_3p=c_\chi\chi(t).
\]

Avec

\[
p(0)=p_0,
\]

la variation des constantes donne

\[
p(t)
=
p_0e^{-\kappa_3t}
+
c_\chi
\int_0^t
e^{-\kappa_3(t-s)}
\chi(s)
\,\mathrm ds.
\]

Lorsque

\[
\kappa_3\neq\lambda_\chi,
\]

on obtient

\[
\boxed{
\begin{aligned}
p(t)
={}&
p_0e^{-\kappa_3t}
+
\frac{c_\chi\chi_\infty}{\kappa_3}
\left(
1-e^{-\kappa_3t}
\right)
\\
&+
c_\chi
\left(
\chi_0-\chi_\infty
\right)
\frac{
e^{-\lambda_\chi t}
-
e^{-\kappa_3t}
}{
\kappa_3-\lambda_\chi
}.
\end{aligned}
}
\]

Dans le cas particulier

\[
\kappa_3=\lambda_\chi,
\]

le dernier quotient doit être remplacé par sa limite :

\[
te^{-\kappa_3t}.
\]

Après un événement,

\[
p_0=r_P\Theta.
\]

Le prochain Kairos se produit au plus petit temps positif
\(\Delta\) tel que

\[
p(\Delta)=\Theta.
\]

Cette équation combine deux exponentielles de taux différents. Elle
est explicite en \(p(t)\), mais transcendante en \(\Delta\).

\subsection{Seuil d'activité événementielle soutenue}
\label{annexe-r2-gk}

En l'absence de nouveau saut, le couple \((\chi,p)\) tend vers

\[
\chi_\infty
=
\frac{gw}{\lambda_\chi},
\]

et

\[
p_\infty
=
\frac{c_\chi}{\kappa_3}\chi_\infty.
\]

Ainsi,

\[
p_\infty
=
\frac{
\kappa_\chi\sqrt2\,gw
}{
\kappa_3\lambda_\chi
}.
\]

Une activité événementielle soutenue est possible lorsque la pression
asymptotique dépasse le seuil :

\[
p_\infty>\Theta.
\]

Le seuil limite est défini par

\[
p_\infty=\Theta.
\]

On obtient

\[
\boxed{
g_K
=
\frac{
\Theta\kappa_3\lambda_\chi
}{
\kappa_\chi\sqrt2\,w
}.
}
\]

Pour les paramètres de référence,

\[
\boxed{
g_K=0{,}03608.
}
\]

À \(g=g_K\), le seuil n'est atteint qu'asymptotiquement dans le modèle
sans perturbation : la période diverge. Pour \(g>g_K\), la pression
finit par franchir \(\Theta\) après chaque relaxation.

Cette formule caractérise l'apparition d'une activité événementielle
soutenue. Elle ne dit rien, à elle seule, sur la stabilité du repos
dans les variables \((x,v)\).

\subsection{Application de retour de \texorpdfstring{$\chi$}{chi}}
\label{annexe-r2-retour-chi}

Notons

\[
\chi_n^+
\]

la charge immédiatement après le \(n\)-ième événement.

Pour une valeur donnée \(\chi_n^+\), le temps
\(\Delta_n\) jusqu'au prochain événement est le plus petit zéro positif
de

\[
p(t;\chi_n^+,r_P\Theta)-\Theta.
\]

La charge juste avant l'événement suivant vaut

\[
\chi_{n+1}^-
=
\chi_\infty
+
\left(
\chi_n^+-\chi_\infty
\right)
e^{-\lambda_\chi\Delta_n}.
\]

Après le saut,

\[
\boxed{
\chi_{n+1}^+
=
(1-r_\chi)
\left[
\chi_\infty
+
\left(
\chi_n^+-\chi_\infty
\right)
e^{-\lambda_\chi\Delta_n}
\right].
}
\]

Cette relation définit une application de retour scalaire

\[
\chi_{n+1}^+
=
\mathcal R_g(\chi_n^+).
\]

Elle est explicite une fois \(\Delta_n\) calculé, mais
\(\Delta_n\) reste la solution d'une équation transcendante.

\subsection{Cycle fixe de l'interface}
\label{annexe-r2-cycle-fixe}

Sur un cycle périodique, la période est constante :

\[
\Delta_n=\Delta.
\]

Posons

\[
E_\chi
=
e^{-\lambda_\chi\Delta}.
\]

Les valeurs pré- et post-événement vérifient

\[
\chi_\ast^+
=
(1-r_\chi)\chi_\ast^-,
\]

et

\[
\chi_\ast^-
=
\chi_\infty
+
\left(
\chi_\ast^+-\chi_\infty
\right)E_\chi.
\]

En éliminant \(\chi_\ast^+\), on obtient

\[
\chi_\ast^-
\left[
1-(1-r_\chi)E_\chi
\right]
=
\chi_\infty
\left(
1-E_\chi
\right).
\]

Ainsi,

\[
\boxed{
\chi_\ast^-(\Delta)
=
\chi_\infty
\frac{
1-e^{-\lambda_\chi\Delta}
}{
1-(1-r_\chi)e^{-\lambda_\chi\Delta}
}.
}
\]

Et

\[
\boxed{
\chi_\ast^+(\Delta)
=
(1-r_\chi)\chi_\ast^-(\Delta).
}
\]

La période est déterminée par l'équation

\[
\boxed{
p
\left(
\Delta;
\chi_\ast^+(\Delta),
r_P\Theta
\right)
=
\Theta.
}
\]

Il s'agit d'une équation scalaire transcendante en \(\Delta\).

Pour les paramètres étudiés et \(g>g_K\), la résolution numérique
donne un cycle attractif de l'interface. L'existence du seuil \(g_K\)
et les formules précédentes sont analytiques ; l'unicité et
l'attractivité globales du cycle pour toute valeur admissible des
paramètres ne sont pas démontrées ici.

\subsection{Application de retour de la trace}
\label{annexe-r2-retour-trace}

Entre deux événements,

\[
\theta(t)
=
\theta_n^+e^{-\beta t}.
\]

Juste avant l'événement suivant,

\[
\theta_{n+1}^-
=
e^{-\beta\Delta}\theta_n^+.
\]

Le quantum ajoute

\[
\kappa_Q\chi_\ast^-.
\]

La carte de la trace est donc

\[
\boxed{
\theta_{n+1}^+
=
e^{-\beta\Delta}\theta_n^+
+
\kappa_Q\chi_\ast^-.
}
\]

Il s'agit d'une application affine contractante, car

\[
0<e^{-\beta\Delta}<1.
\]

Son point fixe est

\[
\boxed{
\theta_\infty^+
=
\frac{
\kappa_Q\chi_\ast^-
}{
1-e^{-\beta\Delta}
}.
}
\]

La valeur juste avant l'événement est

\[
\boxed{
\theta_\infty^-
=
e^{-\beta\Delta}
\theta_\infty^+.
}
\]

\subsection{Moyenne temporelle exacte de la trace}
\label{annexe-r2-moyenne-trace}

Sur le cycle fixe,

\[
\theta_\ast(t)
=
\theta_\infty^+
e^{-\beta t},
\qquad
0\leq t<\Delta.
\]

Sa moyenne temporelle est

\[
\bar\theta
=
\frac{1}{\Delta}
\int_0^\Delta
\theta_\infty^+
e^{-\beta t}
\,\mathrm dt.
\]

Donc

\[
\bar\theta
=
\frac{
\theta_\infty^+
\left(
1-e^{-\beta\Delta}
\right)
}{
\beta\Delta
}.
\]

En utilisant la formule du point fixe,

\[
\boxed{
\bar\theta
=
\frac{
\kappa_Q\chi_\ast^-
}{
\beta\Delta
}.
}
\]

La quantité

\[
\frac{\kappa_Q\chi_\ast^-}{\Delta}
\]

est le débit moyen de trace livré par les événements. La trace moyenne
est ce débit divisé par le taux d'oubli \(\beta\).

\begin{quote}
\resultheading{Proposition}{prop:E-trace-moyenne-R2}{}
Dans le cycle événementiel stationnaire de R2,
\[
\text{trace moyenne}
=
\frac{
\text{quantum par événement}
}{
\text{période}\times\text{taux d'oubli}
}.
\]
\end{quote}

Cette identité est exacte une fois \(\Delta\) connu.

\subsection{Vérification de la loi moyenne}
\label{annexe-r2-verification-moyenne}

Pour quatre intensités testées, les valeurs prédites et mesurées de
\(\bar\theta\) sont :

\[
\begin{array}{c|cc}
 & \text{prédite} & \text{mesurée}\\
\hline
1&0{,}1293&0{,}1288\\
2&0{,}2886&0{,}2885\\
3&0{,}3953&0{,}3952\\
4&0{,}7216&0{,}7216
\end{array}
\]

Les écarts observés viennent de l'intégration et de la détection
numériques des événements. La relation analytique elle-même ne
contient aucun paramètre ajusté. Pour une reproduction complète, le
rapport numérique associé doit exporter, pour chaque ligne, le couple
\((\Delta(g),\chi_\ast^-(g))\) utilisé dans
\(\bar\theta=\kappa_Q\chi_\ast^-/(\beta\Delta)\). Ces valeurs
intermédiaires ne sont pas reconstruites à partir des seules valeurs
arrondies du tableau.

\subsection{Dynamique linéaire de \texorpdfstring{$(x,v)$}{(x,v)} sur le cycle}
\label{annexe-r2-dynamique-lineaire}

Le cycle de l'interface crée une trace en dent de scie. Sur la
diagonale,

\[
\tau(t)
=
(\theta_\ast(t),\theta_\ast(t)),
\]

donc

\[
r_\ast(t)
=
\|\tau(t)\|
=
\sqrt2\,\theta_\ast(t).
\]

La linéarisation de la dynamique de l'état autour du repos est

\[
\frac{\mathrm d}{\mathrm dt}
\begin{pmatrix}
\delta x\\
\delta v
\end{pmatrix}
=
A_g(t)
\begin{pmatrix}
\delta x\\
\delta v
\end{pmatrix},
\]

avec

\[
\boxed{
A_g(t)
=
\begin{pmatrix}
0&1\\
\zeta-m(r_\ast(t))
&
-\Gamma(r_\ast(t))
\end{pmatrix}.
}
\]

La matrice \(A_g(t)\) est périodique de période \(\Delta(g)\).

Notons \(\Phi_g(t)\) la matrice fondamentale définie par

\[
\dot\Phi_g(t)
=
A_g(t)\Phi_g(t),
\qquad
\Phi_g(0)=I_2.
\]

Lors du Kairos, la position est inchangée et la vitesse est relaxée :

\[
\delta x^+=\delta x^-,
\qquad
\delta v^+=r_V\delta v^-.
\]

La matrice de saut linéaire est donc

\[
R_V
=
\begin{pmatrix}
1&0\\
0&r_V
\end{pmatrix}.
\]

La matrice de monodromie, d'un instant immédiatement postérieur à un
événement au suivant, est

\[
\boxed{
\mathcal M_g
=
R_V\Phi_g(\Delta(g)).
}
\]

Le repos périodique est linéairement stable dans cette réduction si

\[
\rho_{\mathrm{sp}}(\mathcal M_g)<1,
\]

où \(\rho_{\mathrm{sp}}\) désigne le rayon spectral.

Il perd sa stabilité lorsque le multiplicateur dominant traverse le
cercle unité :

\[
\boxed{
\rho_{\mathrm{sp}}(\mathcal M_{g_F})=1.
}
\]

La période \(\Delta(g)\) et la matrice fondamentale doivent être
calculées numériquement. Le multiplicateur de Floquet n'est donc pas
disponible sous forme fermée.

\subsection{Portée de la monodromie utilisée}
\label{annexe-r2-portee-monodromie}

La formule précédente est la monodromie du cycle réduit utilisé dans
le Tome II :

\begin{itemize}
\item
  l'horaire des événements est déterminé par le sous-système
  \((\chi,p)\) au repos ;
\item
  les variables \((x,v)\) évoluent linéairement le long de cet horaire ;
\item
  la relaxation de \(v\) est incluse par \(R_V\).
\end{itemize}

Dans le noyau complet, la pression contient aussi des termes en
\(\|x\|\) et \(\|v\|\). Ces normes ne sont pas différentiables au
repos. Une analyse de saltation complète demanderait soit une
régularisation, soit une étude directionnelle de cette dépendance des
instants d'événement.

Le seuil calculé ci-dessous a donc le statut de seuil de stabilité du
cycle R2 réduit, celui effectivement utilisé dans les expériences du
Tome II.

\subsection{Seuil de Floquet}
\label{annexe-r2-gf}

Le calcul numérique du multiplicateur dominant donne

\[
\boxed{
g_F=0{,}0763.
}
\]

Ce seuil est l'analogue local du seuil de stabilité du repos dans un
système à coefficients périodiques.

Il est distinct de \(g_K\) :

\[
g_K<g_F.
\]

Pour

\[
g_K<g<g_F,
\]

l'interface produit déjà des événements soutenus, mais le repos reste
stable dans les variables \((x,v)\).

\subsection{Approximation quasi-statique}
\label{annexe-r2-quasi-statique}

Définissons le seuil scalaire par composante

\[
\theta_{\mathrm{ign}}
=
\frac{\tau_{\mathrm{ign}}}{\sqrt2}.
\]

Une approximation simple remplace la trace périodique par sa moyenne
et impose

\[
\bar\theta(g)
=
\theta_{\mathrm{ign}}.
\]

Le seuil obtenu est

\[
\boxed{
g_F^{qs}=0{,}0760.
}
\]

L'écart relatif avec le calcul de Floquet est d'environ

\[
0{,}4\,\%.
\]

Cette proximité s'explique par la séparation des temps :

\[
\Delta(g)
\simeq3\text{ à }4,
\]

alors que

\[
\frac{1}{\beta}
\simeq8{,}3.
\]

La trace varie peu à l'échelle d'un cycle. La moyenne conserve donc
l'information locale essentielle.

Le symbole \(g_F^{qs}\) ne désigne pas un cinquième seuil. Il s'agit
d'une approximation de \(g_F\).

\subsection{Validation, instabilité et appropriation}
\label{annexe-r2-validation-appropriation}

Trois événements logiques doivent être séparés.

\subsubsection{Activité événementielle}

\[
g>g_K
\]

signifie que l'interface peut produire indéfiniment des Kairos sous
contact soutenu.

\subsubsection{Perte de stabilité du repos}

\[
g>g_F
\]

signifie que le cycle périodique du repos est localement instable dans
la réduction de Floquet.

\subsubsection{Appropriation}

L'appropriation exige que la trajectoire franchisse une séparatrice et
rejoigne, après découplage, un autre attracteur autonome.

Pour la perturbation initiale et le protocole de référence,

\[
\boxed{
g_{\mathrm{app}}^{R2}
(x_0;\mathcal P)
=
0{,}1632.
}
\]

Cette valeur dépend :

\begin{itemize}
\item
  de la perturbation initiale \(x_0\) ;
\item
  du protocole \(\mathcal P\) ;
\item
  de la phase du cycle événementiel ;
\item
  de la structure des bassins.
\end{itemize}

Elle n'est donc ni \(g_K\), ni \(g_F\), ni un seuil structurel
universel.

À \(g=0{,}10\), le protocole de référence produit

\[
271
\]

Kairos et une trace maximale supérieure au seuil d'instabilité, sans
appropriation. Le système est capturé par une configuration résistante
sous charge pulsée.

Ainsi,

\[
\boxed{
\text{validation événementielle}
\neq
\text{appropriation}.
}
\]

\subsection{Résistance pulsée}
\label{annexe-r2-resistance-pulsee}

Dans le modèle continu R0, la branche résistante perd sa stabilité à

\[
g_H^0=0{,}1285.
\]

Sous R2, une rampe lente initialisée sur l'attracteur résistant pulsé
ne montre aucune perte de stabilité jusqu'à

\[
g=0{,}30.
\]

La conclusion correcte est donc

\[
\boxed{
g_H^{R2}>0{,}30
}
\]

au sens d'une borne expérimentale.

Cette écriture ne signifie pas qu'un seuil de Hopf a été localisé
au-delà de \(0{,}30\). Plusieurs possibilités restent ouvertes :

\begin{itemize}
\item
  un seuil plus lointain ;
\item
  une autre bifurcation ;
\item
  une résistance stable pour toute l'étendue testable ;
\item
  un changement d'attracteur suivi par la rampe.
\end{itemize}

\subsection{Hiérarchie des quatre seuils}
\label{annexe-r2-hierarchie}

Pour le protocole et les paramètres étudiés,

\[
\boxed{
g_K
<
g_F
<
g_{\mathrm{app}}^{R2}
<
g_H^{R2}.
}
\]

Les quatre objets sont :

\[
\begin{array}{c|c|l}
\text{symbole}
&
\text{valeur}
&
\text{nature}
\\
\hline
g_K
&
0{,}03608
&
\text{activité événementielle soutenue}
\\
g_F
&
0{,}0763
&
\text{perte de stabilité du repos périodique}
\\
g_{\mathrm{app}}^{R2}(x_0;\mathcal P)
&
0{,}1632
&
\text{franchissement de bassin et persistance}
\\
g_H^{R2}
&
>0{,}30
&
\text{aucune rupture résistante observée avant cette borne}
\end{array}
\]

Le seuil quasi-statique

\[
g_F^{qs}=0{,}0760
\]

est une approximation de \(g_F\), pas un objet supplémentaire.

\section{Comparaison structurale de R1 et R2}
\label{annexe-interfaces-comparaison}

R1 et R2 protègent le système par deux mécanismes distincts.

\subsection{R1 agit sur l'amplitude}

\[
g
\longmapsto
g_{\mathrm{eff}}^\infty(g)
=
\frac{\rho g}{\rho+\eta gw}.
\]

Le couplage est continu, mais plafonné.

\subsection{R2 agit sur la temporalité}

\[
gw
\longmapsto
\left\{
\kappa_Q\chi_n^-
\right\}_{n\geq0}.
\]

Le couplage brut alimente l'interface en continu, mais la trace reçoit
des quanta séparés par une cadence endogène.

La charge continue équivalente est

\[
\boxed{
\text{débit moyen d'écriture}
=
\frac{\kappa_Q\chi_\ast^-}{\Delta}.
}
\]

Ainsi :

\[
\boxed{
\begin{aligned}
\text{R1}
&:
\text{combien peut entrer ?}
\\
\text{R2}
&:
\text{quand et sous quelle forme cela entre-t-il ?}
\end{aligned}
}
\]

\section{Résultats établis et limites}
\label{annexe-interfaces-limites}

\subsection{Résultats analytiques}

Cette annexe établit :

\begin{itemize}
\item
  la solution exacte de R1 ;
\item
  le couplage effectif et son plafond ;
\item
  la formule de transport des seuils ;
\item
  la condition d'inaccessibilité d'un seuil ;
\item
  l'absence d'hystérésis dans R1 minimal ;
\item
  les solutions explicites de \(\chi\) et \(p\) entre deux événements ;
\item
  la formule exacte de \(g_K\) ;
\item
  l'application de retour de \(\chi\) ;
\item
  la caractérisation transcendante de \(\Delta\) ;
\item
  la carte affine de la trace ;
\item
  la formule exacte de sa moyenne temporelle ;
\item
  la construction de la matrice de monodromie réduite.
\end{itemize}

\subsection{Résultats analytiques complétés numériquement}

Les valeurs suivantes demandent une résolution ou une intégration
numérique :

\[
g_F=0{,}0763,
\]

\[
g_F^{qs}=0{,}0760,
\]

les périodes \(\Delta(g)\), et les valeurs de
\(T^\ast\) sous R1.

\subsection{Faits de bassin ou de protocole}

Les quantités suivantes ne sont pas des seuils universels :

\[
g_{\mathrm{app}}^{R2}
(x_0;\mathcal P)
=
0{,}1632,
\]

et

\[
g_H^{R2}>0{,}30.
\]

La première dépend du protocole et de la condition initiale. La seconde
est seulement une borne issue d'une rampe numérique.

\subsection{Questions laissées à l'annexe numérique}

Ne sont pas développés ici :

\begin{itemize}
\item
  les bandes de sommation selon le délai ;
\item
  le prédicteur de phase ;
\item
  les contrôles à charge nette comparable ;
\item
  l'amorçage scalaire ;
\item
  la mémoire vectorielle et la compatibilité directionnelle.
\end{itemize}

Ces résultats demandent surtout une description précise des protocoles
et des validations numériques. Ils relèvent donc de l'annexe suivante.
